\documentclass[12pt,a4paper]{article}

% ----------------------------------------------------------------
% Язык и шрифты: вариант для pdfLaTeX
% ----------------------------------------------------------------
\usepackage[T2A]{fontenc}
\usepackage[utf8]{inputenc}
\usepackage[english,russian]{babel}

\usepackage{cmap}
\usepackage{lmodern}

% ----------------------------------------------------------------
% Математика и оформление
% ----------------------------------------------------------------
\usepackage{amsmath,amssymb,amsthm,mathtools}
\usepackage{bm}
\usepackage{booktabs}
\usepackage{array}
\usepackage{enumitem}
\usepackage{geometry}
\usepackage[expansion=false]{microtype}
\usepackage{xcolor}
\usepackage{hyperref}
\usepackage{tikz}
\usetikzlibrary{positioning,arrows.meta,fit,calc}

\geometry{left=2cm,right=2cm,top=2cm,bottom=2cm}
\hypersetup{
    colorlinks=true,
    linkcolor=blue!45!black,
    urlcolor=blue!55!black,
    citecolor=blue!45!black
}
\allowdisplaybreaks
\setlength{\parindent}{1.25cm}
\setlength{\parskip}{0.15em}

% ----------------------------------------------------------------
% Теоремы и определения
% ----------------------------------------------------------------
\newtheorem{theorem}{Теорема}[section]
\newtheorem{lemma}[theorem]{Лемма}
\newtheorem{proposition}[theorem]{Утверждение}
\newtheorem{corollary}[theorem]{Следствие}

\theoremstyle{definition}
\newtheorem{definition}[theorem]{Определение}
\newtheorem{example}[theorem]{Пример}
\newtheorem{remark}[theorem]{Замечание}

\renewcommand{\proofname}{Доказательство}

% ----------------------------------------------------------------
% Обозначения
% ----------------------------------------------------------------
\newcommand{\R}{\mathbb{R}}
\newcommand{\E}{\mathbb{E}}
\newcommand{\N}{\mathcal{N}}
\newcommand{\KL}{\operatorname{KL}}
\newcommand{\softmax}{\operatorname{softmax}}
\newcommand{\diag}{\operatorname{diag}}
\newcommand{\Var}{\operatorname{Var}}
\newcommand{\Cov}{\operatorname{Cov}}
\newcommand{\tr}{\operatorname{tr}}
\newcommand{\given}{\,\middle|\,}
\newcommand{\dd}{\mathop{}\!\mathrm{d}}
\newcommand{\eps}{\varepsilon}
\newcommand{\1}{\mathbf{1}}

% ----------------------------------------------------------------
% Окружение для алгоритмов
% ----------------------------------------------------------------
\newenvironment{algorithmbox}[1]
{
    \begin{center}
    \begin{minipage}{0.95\textwidth}
    \hrule
    \medskip
    \textbf{Алгоритм: #1}
    \begin{enumerate}[leftmargin=*,label=\arabic*.]
}
{
    \end{enumerate}
    \medskip
    \hrule
    \end{minipage}
    \end{center}
}

\title{Билет №99 \\ \small Трансформеры. Диффузионные модели. (ИТМО)}
\author{Сериков Владислав Александрович}
\date{19 июля 2026}

\begin{document}

\maketitle
\tableofcontents
\newpage

% ================================================================
\section{Основные обозначения и постановка задач}
% ================================================================

Вектор обозначается строчной полужирной буквой, например
$\bm{x}\in\R^d$, а матрица --- прописной полужирной буквой
$\bm{X}\in\R^{n\times d}$. Индекс строки обычно соответствует позиции
объекта в последовательности, а индекс столбца --- признаку.

Для распределения с плотностью $p$ математическое ожидание обозначается
через $\E_p$. Если распределение очевидно из контекста, нижний индекс
опускается.

Дивергенция Кульбака--Лейблера между плотностями $q$ и $p$ определяется
как
\[
    \KL(q\|p)
    =
    \E_{\bm{x}\sim q}
    \left[
        \log\frac{q(\bm{x})}{p(\bm{x})}
    \right].
\]
Она неотрицательна и равна нулю тогда и только тогда, когда
$q=p$ почти всюду.

Рассматриваются две разные, но часто совместно используемые идеи:
\begin{enumerate}
    \item \textbf{трансформер} --- архитектура нейронной сети, основанная
    на механизме внимания;
    \item \textbf{диффузионная модель} --- вероятностная генеративная
    модель, обучающая обратное преобразование шума в данные.
\end{enumerate}

Трансформер может использоваться как самостоятельная авторегрессионная
модель, как кодировщик признаков или как нейросеть, предсказывающая шум
внутри диффузионной модели.

% ================================================================
\section{Трансформеры}
% ================================================================

\subsection{Последовательности и представления токенов}

Пусть дана последовательность дискретных токенов
\[
    (w_1,w_2,\ldots,w_n),
    \qquad
    w_i\in\{1,\ldots,|\mathcal{V}|\},
\]
где $\mathcal{V}$ --- словарь.

Каждому токену ставится в соответствие векторное представление
\[
    \bm{e}_i=\bm{E}[w_i]\in\R^{d_{\text{model}}},
\]
где
\[
    \bm{E}\in\R^{|\mathcal{V}|\times d_{\text{model}}}
\]
--- обучаемая матрица эмбеддингов.

После объединения строк получается матрица
\[
    \bm{X}
    =
    \begin{bmatrix}
        \bm{e}_1^\top\\
        \vdots\\
        \bm{e}_n^\top
    \end{bmatrix}
    \in\R^{n\times d_{\text{model}}}.
\]

Токенизация может выполняться по словам, символам или подсловам.
Подсловная токенизация позволяет представлять редкие и ранее не
встречавшиеся слова через более короткие элементы.

\subsection{Механизм внимания}

\subsubsection{Запросы, ключи и значения}

\begin{definition}
Пусть
\[
    \bm{X}_Q\in\R^{n_q\times d_{\text{model}}},
    \qquad
    \bm{X}_{KV}\in\R^{n_k\times d_{\text{model}}}.
\]
Матрицы запросов, ключей и значений определяются как
\[
    \bm{Q}=\bm{X}_Q\bm{W}^Q,\qquad
    \bm{K}=\bm{X}_{KV}\bm{W}^K,\qquad
    \bm{V}=\bm{X}_{KV}\bm{W}^V,
\]
где
\[
    \bm{W}^Q,\bm{W}^K\in\R^{d_{\text{model}}\times d_k},
    \qquad
    \bm{W}^V\in\R^{d_{\text{model}}\times d_v}.
\]
Масштабированным скалярным вниманием называется отображение
\[
    \operatorname{Attention}(\bm{Q},\bm{K},\bm{V})
    =
    \softmax\left(
        \frac{\bm{Q}\bm{K}^{\top}}{\sqrt{d_k}}+\bm{M}
    \right)\bm{V}.
\]
Здесь softmax применяется независимо к каждой строке, а
$\bm{M}$ --- аддитивная маска.
\end{definition}

Для матрицы $\bm{S}\in\R^{n_q\times n_k}$ построчный softmax имеет вид
\[
    [\softmax(\bm{S})]_{ij}
    =
    \frac{\exp(S_{ij})}
         {\sum_{\ell=1}^{n_k}\exp(S_{i\ell})}.
\]

Элемент
\[
    S_{ij}=\frac{\langle \bm{q}_i,\bm{k}_j\rangle}{\sqrt{d_k}}
\]
характеризует степень соответствия запроса в позиции $i$ ключу в
позиции $j$. После нормировки веса используются для вычисления
взвешенной суммы значений:
\[
    \bm{o}_i=\sum_{j=1}^{n_k}A_{ij}\bm{v}_j,
    \qquad
    \bm{A}=\softmax(\bm{S}+\bm{M}).
\]

\subsubsection{Зачем нужен множитель \(1/\sqrt{d_k}\)}

Пусть компоненты $\bm{q}$ и $\bm{k}$ независимы, имеют нулевое среднее
и единичную дисперсию. Тогда
\[
    \bm{q}^{\top}\bm{k}
    =
    \sum_{\ell=1}^{d_k}q_\ell k_\ell.
\]
При сделанных предположениях
\[
    \Var(q_\ell k_\ell)=1,
\]
поэтому
\[
    \Var(\bm{q}^{\top}\bm{k})=d_k.
\]
Следовательно, типичный масштаб скалярного произведения растёт как
$\sqrt{d_k}$. Деление на $\sqrt{d_k}$ стабилизирует дисперсию логитов
внимания. Без масштабирования softmax при больших $d_k$ может попадать
в область насыщения, где его градиенты малы.

\subsubsection{Маски внимания}

Аддитивная маска определяется как
\[
    M_{ij}
    =
    \begin{cases}
        0, & \text{если внимание из позиции \(i\) в позицию \(j\)
        разрешено},\\
        -\infty, & \text{если оно запрещено}.
    \end{cases}
\]
После softmax запрещённые элементы получают нулевой вес.

Основные виды масок:
\begin{enumerate}
    \item \textbf{padding mask}: запрещает обращение к позициям
    искусственного заполнения;
    \item \textbf{causal mask}: запрещает позиции $i$ использовать
    будущие токены $j>i$:
    \[
        M_{ij}
        =
        \begin{cases}
            0, & j\leq i,\\
            -\infty, & j>i.
        \end{cases}
    \]
    \item \textbf{структурные маски}: задают локальное, блочное или
    разреженное внимание.
\end{enumerate}

\subsubsection{Геометрическое свойство внимания}

\begin{proposition}[Выход внимания как выпуклая комбинация]
Пусть все немаскированные логиты конечны. Тогда каждая строка выхода
\[
    \bm{O}
    =
    \softmax\left(
        \frac{\bm{Q}\bm{K}^{\top}}{\sqrt{d_k}}+\bm{M}
    \right)\bm{V}
\]
является выпуклой комбинацией строк матрицы $\bm{V}$, доступных данной
позиции.
\end{proposition}

\begin{proof}
Обозначим матрицу весов через
\[
    \bm{A}
    =
    \softmax\left(
        \frac{\bm{Q}\bm{K}^{\top}}{\sqrt{d_k}}+\bm{M}
    \right).
\]
Для каждого разрешённого элемента
\[
    A_{ij}
    =
    \frac{\exp(S_{ij})}
         {\sum_{\ell:\,M_{i\ell}=0}\exp(S_{i\ell})}>0.
\]
Для запрещённых элементов после применения предельного значения
$\exp(-\infty)=0$ имеем $A_{ij}=0$.

По определению softmax сумма элементов каждой строки равна единице:
\[
    \sum_{j=1}^{n_k}A_{ij}=1.
\]
Строка выхода равна
\[
    \bm{o}_i=\sum_{j=1}^{n_k}A_{ij}\bm{v}_j.
\]
Таким образом, коэффициенты $A_{ij}$ неотрицательны и в сумме дают
единицу. Это и означает, что $\bm{o}_i$ является выпуклой комбинацией
доступных строк $\bm{V}$.
\end{proof}

\subsubsection{Минимальный численный пример}

Пусть
\[
    \bm{X}
    =
    \begin{bmatrix}
        1&0\\
        0&1
    \end{bmatrix},
    \qquad
    \bm{W}^Q=\bm{W}^K=\bm{W}^V=\bm{I}_2.
\]
Тогда
\[
    \bm{Q}=\bm{K}=\bm{V}=\bm{X},
    \qquad
    \frac{\bm{Q}\bm{K}^{\top}}{\sqrt{2}}
    =
    \begin{bmatrix}
        1/\sqrt{2}&0\\
        0&1/\sqrt{2}
    \end{bmatrix}.
\]
Обозначим
\[
    a
    =
    \frac{e^{1/\sqrt{2}}}
         {e^{1/\sqrt{2}}+1}
    \approx 0.6698.
\]
Тогда
\[
    \bm{A}
    \approx
    \begin{bmatrix}
        0.6698&0.3302\\
        0.3302&0.6698
    \end{bmatrix},
\]
и
\[
    \bm{O}=\bm{A}\bm{V}
    \approx
    \begin{bmatrix}
        0.6698&0.3302\\
        0.3302&0.6698
    \end{bmatrix}.
\]

С причинной маской
\[
    \bm{M}
    =
    \begin{bmatrix}
        0&-\infty\\
        0&0
    \end{bmatrix}
\]
получаем
\[
    \bm{A}_{\text{causal}}
    \approx
    \begin{bmatrix}
        1&0\\
        0.3302&0.6698
    \end{bmatrix}.
\]
Первая позиция не использует вторую, поскольку она является будущей.

\subsection{Self-attention и cross-attention}

\begin{definition}
В \textbf{self-attention} запросы, ключи и значения строятся из одной
и той же последовательности:
\[
    \bm{X}_Q=\bm{X}_{KV}=\bm{X}.
\]
\end{definition}

Self-attention позволяет каждой позиции агрегировать информацию из
остальных позиций той же последовательности.

\begin{definition}
В \textbf{cross-attention} запросы строятся из одной
последовательности, а ключи и значения --- из другой:
\[
    \bm{Q}=\bm{X}_{\text{target}}\bm{W}^Q,
    \qquad
    \bm{K}=\bm{X}_{\text{source}}\bm{W}^K,
    \qquad
    \bm{V}=\bm{X}_{\text{source}}\bm{W}^V.
\]
\end{definition}

Например, в машинном переводе запросы декодера обращаются к выходам
кодировщика, содержащим представление исходного предложения.

\subsection{Многоголовое внимание}

Одна голова внимания использует один набор проекций и одну матрицу
сходства. Многоголовое внимание вычисляет несколько механизмов
внимания параллельно:
\[
    \operatorname{head}_r
    =
    \operatorname{Attention}
    \left(
        \bm{X}_Q\bm{W}_r^Q,\,
        \bm{X}_{KV}\bm{W}_r^K,\,
        \bm{X}_{KV}\bm{W}_r^V
    \right),
\]
\[
    \operatorname{MHA}(\bm{X}_Q,\bm{X}_{KV})
    =
    \operatorname{Concat}
    (\operatorname{head}_1,\ldots,\operatorname{head}_h)\bm{W}^O.
\]
Обычно выбирают
\[
    d_k=d_v=\frac{d_{\text{model}}}{h}.
\]
Тогда конкатенация голов снова имеет размерность
$d_{\text{model}}$.

Разные головы могут специализироваться на различных типах
зависимостей: локальных, синтаксических, позиционных или семантических.
Такая интерпретация не гарантируется формально, но часто наблюдается
эмпирически.

\subsection{Позиционная информация}

\subsubsection{Необходимость позиционного кодирования}

Без позиционной информации self-attention не различает порядок
элементов в последовательности.

\begin{theorem}[Перестановочная эквивариантность self-attention]
Пусть
\[
    \operatorname{SA}(\bm{X})
    =
    \softmax\left(
        \frac{\bm{X}\bm{W}^Q
        (\bm{X}\bm{W}^K)^{\top}}{\sqrt{d_k}}
    \right)
    \bm{X}\bm{W}^V
\]
--- self-attention без позиционных признаков и без маски. Для любой
перестановочной матрицы $\bm{P}\in\R^{n\times n}$ выполняется
\[
    \operatorname{SA}(\bm{P}\bm{X})
    =
    \bm{P}\operatorname{SA}(\bm{X}).
\]
Иными словами, перестановка входных позиций приводит только к такой
же перестановке выходов.
\end{theorem}

\begin{proof}
Положим
\[
    \bm{Q}=\bm{X}\bm{W}^Q,\qquad
    \bm{K}=\bm{X}\bm{W}^K,\qquad
    \bm{V}=\bm{X}\bm{W}^V.
\]
После перестановки строк входа:
\[
    \bm{Q}'=\bm{P}\bm{Q},\qquad
    \bm{K}'=\bm{P}\bm{K},\qquad
    \bm{V}'=\bm{P}\bm{V}.
\]
Матрица логитов преобразуется следующим образом:
\[
    \bm{Q}'(\bm{K}')^\top
    =
    \bm{P}\bm{Q}\bm{K}^\top\bm{P}^\top.
\]

Докажем вспомогательное равенство
\[
    \softmax(\bm{P}\bm{S}\bm{P}^\top)
    =
    \bm{P}\softmax(\bm{S})\bm{P}^\top.
\]
Пусть перестановка, задаваемая $\bm{P}$, обозначена через $\pi$.
Элемент левой части имеет вид
\[
    \frac{
        \exp(S_{\pi(i),\pi(j)})
    }{
        \sum_{k=1}^{n}\exp(S_{\pi(i),\pi(k)})
    }.
\]
Поскольку $\pi$ является биекцией, множество индексов
$\{\pi(k):k=1,\ldots,n\}$ совпадает с $\{1,\ldots,n\}$. Поэтому
знаменатель равен
\[
    \sum_{\ell=1}^{n}\exp(S_{\pi(i),\ell}),
\]
а весь элемент совпадает с элементом
$(i,j)$ матрицы
$\bm{P}\softmax(\bm{S})\bm{P}^\top$.

Следовательно,
\begin{align*}
    \operatorname{SA}(\bm{P}\bm{X})
    &=
    \softmax\left(
        \bm{P}
        \frac{\bm{Q}\bm{K}^\top}{\sqrt{d_k}}
        \bm{P}^\top
    \right)\bm{P}\bm{V}\\
    &=
    \bm{P}
    \softmax\left(
        \frac{\bm{Q}\bm{K}^\top}{\sqrt{d_k}}
    \right)
    \bm{P}^\top\bm{P}\bm{V}\\
    &=
    \bm{P}
    \softmax\left(
        \frac{\bm{Q}\bm{K}^\top}{\sqrt{d_k}}
    \right)\bm{V}\\
    &=
    \bm{P}\operatorname{SA}(\bm{X}).
\end{align*}
Теорема доказана.
\end{proof}

Позиционные признаки нарушают эту симметрию, позволяя сети различать,
например, последовательности
\[
    (A,B,C)
    \quad\text{и}\quad
    (C,B,A).
\]

\subsubsection{Синусоидальное позиционное кодирование}

Для позиции $p$ и индекса признака $i$ можно определить
\[
    \operatorname{PE}(p,2i)
    =
    \sin\left(
        \frac{p}{10000^{2i/d_{\text{model}}}}
    \right),
\]
\[
    \operatorname{PE}(p,2i+1)
    =
    \cos\left(
        \frac{p}{10000^{2i/d_{\text{model}}}}
    \right).
\]
Вход блока равен
\[
    \bm{x}_p=\bm{e}_{w_p}+\operatorname{PE}(p).
\]

Разные пары координат соответствуют разным частотам. Это позволяет
представлять как локальные, так и длинные позиционные зависимости.

\subsubsection{Обучаемые позиционные эмбеддинги}

Для каждой позиции задаётся обучаемый вектор
\[
    \bm{p}_i\in\R^{d_{\text{model}}},
    \qquad
    \bm{x}_i=\bm{e}_{w_i}+\bm{p}_i.
\]
Недостаток такого подхода заключается в том, что позиции за пределами
максимальной длины обучения не имеют непосредственно обученных
представлений.

\subsubsection{Относительные и вращательные позиции}

Относительные методы кодируют не только абсолютные индексы, но и
расстояние $i-j$ между позициями.

В RoPE пары координат запросов и ключей поворачиваются на угол,
зависящий от позиции. Для позиции $m$:
\[
    \bm{q}_m'=\bm{R}_m\bm{q}_m,
    \qquad
    \bm{k}_n'=\bm{R}_n\bm{k}_n,
\]
где $\bm{R}_m$ --- блочно-диагональная ортогональная матрица вращений.

Тогда
\[
    (\bm{q}_m')^\top\bm{k}_n'
    =
    \bm{q}_m^\top\bm{R}_m^\top\bm{R}_n\bm{k}_n
    =
    \bm{q}_m^\top\bm{R}_{n-m}\bm{k}_n.
\]
Таким образом, скалярное произведение явно зависит от относительного
смещения $n-m$.

\subsection{Нормализация, остаточные связи и FFN}

\subsubsection{Layer Normalization}

Для вектора $\bm{x}\in\R^d$ определим
\[
    \mu(\bm{x})=\frac{1}{d}\sum_{j=1}^{d}x_j,
\]
\[
    \sigma^2(\bm{x})
    =
    \frac{1}{d}\sum_{j=1}^{d}(x_j-\mu(\bm{x}))^2.
\]
Тогда
\[
    \operatorname{LN}(\bm{x})
    =
    \bm{\gamma}\odot
    \frac{\bm{x}-\mu(\bm{x})\bm{1}}
         {\sqrt{\sigma^2(\bm{x})+\epsilon}}
    +
    \bm{\beta},
\]
где $\bm{\gamma}$ и $\bm{\beta}$ --- обучаемые параметры.

Нормализация выполняется отдельно для каждого токена по его признакам,
поэтому не зависит от размера мини-пакета.

\subsubsection{Позиционно-независимая полносвязная сеть}

После внимания к каждой позиции независимо применяется одна и та же
сеть:
\[
    \operatorname{FFN}(\bm{x})
    =
    \phi(\bm{x}\bm{W}_1+\bm{b}_1)\bm{W}_2+\bm{b}_2,
\]
где
\[
    \bm{W}_1\in\R^{d_{\text{model}}\times d_{\text{ff}}},
    \qquad
    \bm{W}_2\in\R^{d_{\text{ff}}\times d_{\text{model}}},
\]
а $\phi$ --- нелинейность, например ReLU, GELU или SwiGLU-подобная
функция.

Внимание смешивает информацию между позициями, а FFN выполняет
нелинейное преобразование признаков внутри каждой позиции.

\subsubsection{Pre-LN и Post-LN}

В варианте Post-LN подслой вычисляется как
\[
    \bm{Y}
    =
    \operatorname{LN}
    \bigl(
        \bm{X}+\operatorname{Sublayer}(\bm{X})
    \bigr).
\]

В варианте Pre-LN:
\[
    \bm{Y}
    =
    \bm{X}
    +
    \operatorname{Sublayer}
    \bigl(
        \operatorname{LN}(\bm{X})
    \bigr).
\]

Для Pre-LN блок кодировщика можно записать так:
\[
    \bm{U}
    =
    \bm{X}
    +
    \operatorname{MHA}
    \bigl(
        \operatorname{LN}(\bm{X}),
        \operatorname{LN}(\bm{X})
    \bigr),
\]
\[
    \bm{Y}
    =
    \bm{U}
    +
    \operatorname{FFN}
    \bigl(
        \operatorname{LN}(\bm{U})
    \bigr).
\]
На практике также используются dropout и другие виды регуляризации.

\subsection{Архитектуры трансформеров}

\subsubsection{Кодировщик}

Кодировщик состоит из повторяющихся блоков
\[
    \text{self-attention}
    \longrightarrow
    \text{FFN}.
\]
Self-attention обычно является двунаправленным: каждая позиция может
обращаться ко всем непустым позициям.

Кодировщик применим к:
\begin{itemize}
    \item классификации текста;
    \item извлечению признаков;
    \item разметке последовательностей;
    \item поиску и построению эмбеддингов;
    \item обработке изображений, аудио и других модальностей.
\end{itemize}

\subsubsection{Декодировщик}

Авторегрессионный декодировщик содержит причинное self-attention:
\[
    p_\theta(w_{1:n})
    =
    \prod_{t=1}^{n}
    p_\theta(w_t\mid w_{<t}).
\]
Маска гарантирует, что представление позиции $t$ не зависит от токенов
$w_{t+1},\ldots,w_n$.

\subsubsection{Кодировщик--декодировщик}

Классическая архитектура преобразования последовательностей состоит
из:
\begin{enumerate}
    \item кодировщика исходной последовательности;
    \item причинного self-attention в декодировщике;
    \item cross-attention декодировщика к выходам кодировщика;
    \item позиционной FFN.
\end{enumerate}

\begin{figure}[h]
\centering
\begin{tikzpicture}[
    box/.style={
        draw,
        rounded corners,
        minimum width=4.2cm,
        minimum height=0.8cm,
        align=center,
        fill=blue!6
    },
    arrow/.style={-{Latex[length=2.5mm]},thick},
    node distance=0.65cm and 1.5cm
]
\node[box] (src) {Исходные токены};
\node[box,below=of src] (srcemb) {Эмбеддинги и позиции};
\node[box,below=of srcemb] (enc) {Блоки кодировщика\\Self-attention + FFN};

\node[box,right=of src] (tgt) {Сдвинутые целевые токены};
\node[box,below=of tgt] (tgtemb) {Эмбеддинги и позиции};
\node[box,below=of tgtemb] (decself) {Причинное self-attention};
\node[box,below=of decself] (cross) {Cross-attention};
\node[box,below=of cross] (ffn) {FFN};
\node[box,below=of ffn] (out) {Логиты следующего токена};

\draw[arrow] (src)--(srcemb);
\draw[arrow] (srcemb)--(enc);
\draw[arrow] (tgt)--(tgtemb);
\draw[arrow] (tgtemb)--(decself);
\draw[arrow] (decself)--(cross);
\draw[arrow] (cross)--(ffn);
\draw[arrow] (ffn)--(out);
\draw[arrow] (enc.east)--(cross.west);
\end{tikzpicture}
\caption{Схема трансформера типа кодировщик--декодировщик.}
\end{figure}

\subsection{Функции потерь}

\subsubsection{Авторегрессионное моделирование}

Для последовательности $\bm{w}=(w_1,\ldots,w_n)$ минимизируется
отрицательное логарифмическое правдоподобие
\[
    \mathcal{L}_{\text{AR}}(\theta)
    =
    -\sum_{t=1}^{n}
    \log p_\theta(w_t\mid w_{<t}).
\]
При наличии контекста $\bm{c}$:
\[
    \mathcal{L}_{\text{cond}}(\theta)
    =
    -\sum_{t=1}^{n}
    \log p_\theta(w_t\mid w_{<t},\bm{c}).
\]

Во время обучения все истинные предыдущие токены доступны одновременно.
Такой режим называется \textbf{teacher forcing}. Причинная маска
сохраняет корректную факторизацию вероятности.

\subsubsection{Маскированное моделирование}

Выбирается множество позиций $\mathcal{M}$, часть токенов скрывается,
после чего минимизируется
\[
    \mathcal{L}_{\text{MLM}}(\theta)
    =
    -\sum_{i\in\mathcal{M}}
    \log p_\theta
    \left(
        w_i\mid \bm{w}_{\setminus\mathcal{M}}
    \right).
\]
Такая цель естественна для двунаправленных кодировщиков, поскольку
скрытый токен может восстанавливаться по левому и правому контексту.

\subsubsection{Выходное распределение}

Для скрытого состояния $\bm{h}_t$ вычисляются логиты
\[
    \bm{z}_t=\bm{h}_t\bm{W}_{\text{out}}+\bm{b}_{\text{out}}.
\]
Вероятность токена $j$:
\[
    p_\theta(w_t=j\mid w_{<t})
    =
    \frac{\exp(z_{t,j})}
         {\sum_{k=1}^{|\mathcal{V}|}\exp(z_{t,k})}.
\]

Нередко входная матрица эмбеддингов и выходная матрица связываются:
\[
    \bm{W}_{\text{out}}=\bm{E}^{\top}.
\]

\subsection{Алгоритмы работы трансформера}

\begin{algorithmbox}{Прямой проход блока Pre-LN Transformer}
    \item На вход поступает
    $\bm{X}\in\R^{n\times d_{\text{model}}}$.
    \item Вычисляется нормализованный вход
    \[
        \bm{X}_{\text{norm}}=\operatorname{LN}(\bm{X}).
    \]
    \item Для каждой головы строятся запросы, ключи и значения.
    \item Вычисляются логиты
    \[
        \bm{S}_r
        =
        \frac{\bm{Q}_r\bm{K}_r^\top}{\sqrt{d_k}}+\bm{M}.
    \]
    \item К логитам применяется softmax, после чего значения
    агрегируются:
    \[
        \bm{H}_r=\softmax(\bm{S}_r)\bm{V}_r.
    \]
    \item Выходы голов конкатенируются и проецируются:
    \[
        \bm{H}
        =
        \operatorname{Concat}(\bm{H}_1,\ldots,\bm{H}_h)\bm{W}^O.
    \]
    \item Добавляется остаточная связь:
    \[
        \bm{U}=\bm{X}+\bm{H}.
    \]
    \item Выполняется второй подслой:
    \[
        \bm{Y}
        =
        \bm{U}
        +
        \operatorname{FFN}
        \bigl(
            \operatorname{LN}(\bm{U})
        \bigr).
    \]
    \item Матрица $\bm{Y}$ передаётся следующему блоку.
\end{algorithmbox}

\begin{algorithmbox}{Авторегрессионная генерация}
    \item Токенизировать контекст
    $(w_1,\ldots,w_m)$.
    \item Выполнить прямой проход и получить логиты следующего токена.
    \item Преобразовать логиты в распределение с температурой
    $\tau>0$:
    \[
        p_j
        =
        \frac{\exp(z_j/\tau)}
             {\sum_k\exp(z_k/\tau)}.
    \]
    \item Выбрать следующий токен:
    \begin{itemize}
        \item максимум вероятности;
        \item случайная выборка;
        \item top-$k$;
        \item nucleus sampling по минимальному множеству токенов с
        суммарной вероятностью не меньше заданного порога.
    \end{itemize}
    \item Добавить выбранный токен к контексту.
    \item Повторять шаги 2--5 до появления токена конца или достижения
    максимальной длины.
\end{algorithmbox}

\paragraph{Минимальная иллюстрация.}
Пусть после контекста ``нейронная сеть'' модель выдаёт вероятности
\[
    p(\text{обучается})=0.55,\quad
    p(\text{состоит})=0.25,\quad
    p(\text{работает})=0.20.
\]
Жадный алгоритм выберет ``обучается''. При случайном декодировании
возможны все три продолжения с указанными вероятностями.

\subsection{KV-кэш при генерации}

В причинном self-attention ключи и значения уже обработанных токенов
не меняются. Поэтому для каждого слоя можно сохранять
\[
    \bm{K}_{1:t-1},\qquad \bm{V}_{1:t-1}.
\]
Для нового токена вычисляются только
\[
    \bm{q}_t,\qquad \bm{k}_t,\qquad \bm{v}_t,
\]
после чего
\[
    \bm{K}_{1:t}
    =
    \operatorname{Concat}(\bm{K}_{1:t-1},\bm{k}_t),
\]
\[
    \bm{V}_{1:t}
    =
    \operatorname{Concat}(\bm{V}_{1:t-1},\bm{v}_t).
\]

KV-кэш уменьшает количество повторных вычислений, но требует памяти,
линейно растущей с длиной контекста, числом слоёв и размерностью
ключей и значений.

\subsection{Вычислительная сложность}

Для последовательности длины $n$:
\begin{itemize}
    \item построение матрицы $\bm{Q}\bm{K}^{\top}$ требует порядка
    $O(n^2d_k)$ операций;
    \item хранение полной матрицы внимания требует $O(n^2)$ памяти
    на голову;
    \item линейные проекции требуют порядка $O(nd_{\text{model}}^2)$;
    \item FFN требует порядка
    $O(nd_{\text{model}}d_{\text{ff}})$.
\end{itemize}

Основное ограничение стандартного внимания для длинных
последовательностей --- квадратичный рост матрицы внимания.

Возможные способы уменьшения затрат:
\begin{itemize}
    \item локальное внимание;
    \item блочное и разреженное внимание;
    \item низкоранговые аппроксимации;
    \item линейное внимание;
    \item объединение близких токенов;
    \item обработка последовательности по сегментам.
\end{itemize}

\subsection{Основные достоинства и ограничения}

\paragraph{Достоинства.}
\begin{itemize}
    \item Параллельная обработка всех позиций при обучении.
    \item Короткий путь передачи информации между удалёнными
    позициями.
    \item Универсальность относительно модальности данных.
    \item Возможность масштабирования глубины, ширины, числа голов и
    объёма данных.
    \item Естественная поддержка условной генерации через
    cross-attention.
\end{itemize}

\paragraph{Ограничения.}
\begin{itemize}
    \item Квадратичная сложность полного внимания.
    \item Последовательная авторегрессионная генерация.
    \item Большой объём памяти для параметров, активаций и KV-кэша.
    \item Зависимость от токенизации и обучающего распределения.
    \item Отсутствие гарантии фактической достоверности генерируемого
    текста.
\end{itemize}

% ================================================================
\section{Диффузионные модели}
% ================================================================

\subsection{Общая идея}

Пусть наблюдаемые данные имеют неизвестное распределение
\[
    \bm{x}_0\sim q_{\text{data}}(\bm{x}_0).
\]
Прямая диффузия постепенно разрушает структуру данных, добавляя шум:
\[
    \bm{x}_0
    \longrightarrow
    \bm{x}_1
    \longrightarrow
    \cdots
    \longrightarrow
    \bm{x}_T.
\]
При достаточно большом $T$ распределение $\bm{x}_T$ становится близким
к простому априорному распределению, обычно $\N(\bm{0},\bm{I})$.

Затем обучается обратный процесс:
\[
    \bm{x}_T
    \longrightarrow
    \bm{x}_{T-1}
    \longrightarrow
    \cdots
    \longrightarrow
    \bm{x}_0,
\]
который удаляет шум и восстанавливает структуру данных.

\begin{figure}[h]
\centering
\begin{tikzpicture}[
    state/.style={
        draw,
        circle,
        minimum size=1.25cm,
        fill=blue!6
    },
    fwd/.style={-{Latex[length=2.5mm]},thick},
    rev/.style={-{Latex[length=2.5mm]},thick,dashed,red!65!black},
    node distance=1.7cm
]
\node[state] (x0) {$\bm{x}_0$};
\node[state,right=of x0] (x1) {$\bm{x}_1$};
\node[state,right=of x1] (x2) {$\bm{x}_2$};
\node[right=of x2] (dots) {$\cdots$};
\node[state,right=of dots] (xt) {$\bm{x}_T$};

\draw[fwd,bend left=14] (x0) to node[above] {$q$} (x1);
\draw[fwd,bend left=14] (x1) to node[above] {$q$} (x2);
\draw[fwd,bend left=14] (x2) to (dots);
\draw[fwd,bend left=14] (dots) to (xt);

\draw[rev,bend left=14] (xt) to node[below] {$p_\theta$} (dots);
\draw[rev,bend left=14] (dots) to (x2);
\draw[rev,bend left=14] (x2) to node[below] {$p_\theta$} (x1);
\draw[rev,bend left=14] (x1) to node[below] {$p_\theta$} (x0);
\end{tikzpicture}
\caption{Прямой процесс добавляет шум, обратный процесс удаляет его.}
\end{figure}

\subsection{Прямой процесс DDPM}

Выбирается расписание шума
\[
    0<\beta_1<\beta_2<\cdots<\beta_T<1.
\]
Обозначим
\[
    \alpha_t=1-\beta_t,
    \qquad
    \overline{\alpha}_t
    =
    \prod_{s=1}^{t}\alpha_s,
    \qquad
    \overline{\alpha}_0=1.
\]

Прямой марковский процесс задаётся распределениями
\[
    q(\bm{x}_{1:T}\mid\bm{x}_0)
    =
    \prod_{t=1}^{T}
    q(\bm{x}_t\mid\bm{x}_{t-1}),
\]
где
\[
    q(\bm{x}_t\mid\bm{x}_{t-1})
    =
    \N\left(
        \bm{x}_t;
        \sqrt{\alpha_t}\bm{x}_{t-1},
        \beta_t\bm{I}
    \right).
\]
Эквивалентная репараметризация:
\[
    \bm{x}_t
    =
    \sqrt{\alpha_t}\bm{x}_{t-1}
    +
    \sqrt{\beta_t}\bm{z}_t,
    \qquad
    \bm{z}_t\sim\N(\bm{0},\bm{I}).
\]

\subsection{Закрытая формула зашумления}

\begin{theorem}[Распределение \(\bm{x}_t\) при известном \(\bm{x}_0\)]
Для прямого процесса
\[
    q(\bm{x}_t\mid\bm{x}_{t-1})
    =
    \N\left(
        \sqrt{\alpha_t}\bm{x}_{t-1},
        \beta_t\bm{I}
    \right)
\]
выполняется
\[
    q(\bm{x}_t\mid\bm{x}_0)
    =
    \N\left(
        \sqrt{\overline{\alpha}_t}\bm{x}_0,
        (1-\overline{\alpha}_t)\bm{I}
    \right).
\]
Следовательно,
\[
    \bm{x}_t
    =
    \sqrt{\overline{\alpha}_t}\bm{x}_0
    +
    \sqrt{1-\overline{\alpha}_t}\bm{\eps},
    \qquad
    \bm{\eps}\sim\N(\bm{0},\bm{I}).
\]
\end{theorem}

\begin{proof}
Докажем утверждение по индукции.

Для $t=1$:
\[
    \bm{x}_1
    =
    \sqrt{\alpha_1}\bm{x}_0
    +
    \sqrt{\beta_1}\bm{z}_1.
\]
Поскольку
\[
    \overline{\alpha}_1=\alpha_1,
    \qquad
    1-\overline{\alpha}_1
    =
    1-\alpha_1
    =
    \beta_1,
\]
получаем требуемую формулу.

Предположим, что для $t-1$:
\[
    \bm{x}_{t-1}
    =
    \sqrt{\overline{\alpha}_{t-1}}\bm{x}_0
    +
    \sqrt{1-\overline{\alpha}_{t-1}}\bm{\eps}_{t-1},
\]
где
\[
    \bm{\eps}_{t-1}\sim\N(\bm{0},\bm{I}).
\]
Тогда
\begin{align*}
    \bm{x}_t
    &=
    \sqrt{\alpha_t}\bm{x}_{t-1}
    +
    \sqrt{\beta_t}\bm{z}_t\\
    &=
    \sqrt{\alpha_t\overline{\alpha}_{t-1}}\bm{x}_0
    +
    \sqrt{\alpha_t(1-\overline{\alpha}_{t-1})}
    \bm{\eps}_{t-1}
    +
    \sqrt{\beta_t}\bm{z}_t\\
    &=
    \sqrt{\overline{\alpha}_t}\bm{x}_0
    +
    \bm{\eta}_t,
\end{align*}
где
\[
    \bm{\eta}_t
    =
    \sqrt{\alpha_t(1-\overline{\alpha}_{t-1})}
    \bm{\eps}_{t-1}
    +
    \sqrt{\beta_t}\bm{z}_t.
\]
Векторы $\bm{\eps}_{t-1}$ и $\bm{z}_t$ независимы и имеют стандартное
нормальное распределение. Поэтому $\bm{\eta}_t$ также имеет нормальное
распределение с нулевым средним и ковариацией
\begin{align*}
    \Cov(\bm{\eta}_t)
    &=
    \left[
        \alpha_t(1-\overline{\alpha}_{t-1})
        +
        \beta_t
    \right]\bm{I}\\
    &=
    \left[
        \alpha_t
        -
        \alpha_t\overline{\alpha}_{t-1}
        +
        1-\alpha_t
    \right]\bm{I}\\
    &=
    (1-\overline{\alpha}_t)\bm{I}.
\end{align*}
Следовательно,
\[
    \bm{\eta}_t
    \overset{d}{=}
    \sqrt{1-\overline{\alpha}_t}\bm{\eps},
    \qquad
    \bm{\eps}\sim\N(\bm{0},\bm{I}).
\]
Индукционный переход доказан.
\end{proof}

Важное следствие: для обучения не требуется последовательно
выполнять $t$ шагов зашумления. Любой уровень шума можно получить
непосредственно:
\[
    \bm{x}_t
    =
    \sqrt{\overline{\alpha}_t}\bm{x}_0
    +
    \sqrt{1-\overline{\alpha}_t}\bm{\eps}.
\]

Если $\overline{\alpha}_T\approx 0$, то
\[
    q(\bm{x}_T\mid\bm{x}_0)
    \approx
    \N(\bm{0},\bm{I}).
\]

\subsection{Апостериорное распределение прямого процесса}

Для обучения обратного перехода требуется распределение
\[
    q(\bm{x}_{t-1}\mid\bm{x}_t,\bm{x}_0).
\]

\begin{theorem}[Точный апостериорный переход]
Для $t\geq 2$ выполняется
\[
    q(\bm{x}_{t-1}\mid\bm{x}_t,\bm{x}_0)
    =
    \N\left(
        \bm{x}_{t-1};
        \widetilde{\bm{\mu}}_t(\bm{x}_t,\bm{x}_0),
        \widetilde{\beta}_t\bm{I}
    \right),
\]
где
\[
    \widetilde{\beta}_t
    =
    \frac{1-\overline{\alpha}_{t-1}}
         {1-\overline{\alpha}_t}
    \beta_t
\]
и
\[
    \widetilde{\bm{\mu}}_t(\bm{x}_t,\bm{x}_0)
    =
    \frac{
        \sqrt{\overline{\alpha}_{t-1}}\beta_t
    }{
        1-\overline{\alpha}_t
    }\bm{x}_0
    +
    \frac{
        \sqrt{\alpha_t}
        (1-\overline{\alpha}_{t-1})
    }{
        1-\overline{\alpha}_t
    }\bm{x}_t.
\]
\end{theorem}

\begin{proof}
По формуле Байеса и марковскому свойству
\[
    q(\bm{x}_{t-1}\mid\bm{x}_t,\bm{x}_0)
    \propto
    q(\bm{x}_t\mid\bm{x}_{t-1})
    q(\bm{x}_{t-1}\mid\bm{x}_0).
\]
Из определения прямого процесса:
\[
    q(\bm{x}_t\mid\bm{x}_{t-1})
    \propto
    \exp\left(
        -
        \frac{
            \|\bm{x}_t-\sqrt{\alpha_t}\bm{x}_{t-1}\|^2
        }{
            2\beta_t
        }
    \right).
\]
По предыдущей теореме:
\[
    q(\bm{x}_{t-1}\mid\bm{x}_0)
    \propto
    \exp\left(
        -
        \frac{
            \|\bm{x}_{t-1}
            -\sqrt{\overline{\alpha}_{t-1}}\bm{x}_0\|^2
        }{
            2(1-\overline{\alpha}_{t-1})
        }
    \right).
\]

Логарифм произведения с точностью до константы равен
\begin{align*}
    &-
    \frac{1}{2\beta_t}
    \|\bm{x}_t-\sqrt{\alpha_t}\bm{x}_{t-1}\|^2\\
    &\quad
    -
    \frac{1}{2(1-\overline{\alpha}_{t-1})}
    \|\bm{x}_{t-1}
    -\sqrt{\overline{\alpha}_{t-1}}\bm{x}_0\|^2.
\end{align*}

Коэффициент при
$\bm{x}_{t-1}^{\top}\bm{x}_{t-1}$ равен
\[
    \frac{\alpha_t}{\beta_t}
    +
    \frac{1}{1-\overline{\alpha}_{t-1}}.
\]
Приведём к общему знаменателю:
\begin{align*}
    \frac{\alpha_t}{\beta_t}
    +
    \frac{1}{1-\overline{\alpha}_{t-1}}
    &=
    \frac{
        \alpha_t(1-\overline{\alpha}_{t-1})+\beta_t
    }{
        \beta_t(1-\overline{\alpha}_{t-1})
    }\\
    &=
    \frac{
        1-\alpha_t\overline{\alpha}_{t-1}
    }{
        \beta_t(1-\overline{\alpha}_{t-1})
    }\\
    &=
    \frac{
        1-\overline{\alpha}_t
    }{
        \beta_t(1-\overline{\alpha}_{t-1})
    }.
\end{align*}
Следовательно, ковариация равна обратной величине:
\[
    \widetilde{\beta}_t\bm{I}
    =
    \frac{
        \beta_t(1-\overline{\alpha}_{t-1})
    }{
        1-\overline{\alpha}_t
    }\bm{I}.
\]

Линейная по $\bm{x}_{t-1}$ часть имеет коэффициент
\[
    \frac{\sqrt{\alpha_t}}{\beta_t}\bm{x}_t
    +
    \frac{
        \sqrt{\overline{\alpha}_{t-1}}
    }{
        1-\overline{\alpha}_{t-1}
    }\bm{x}_0.
\]
Умножая его на ковариацию, получаем среднее:
\begin{align*}
    \widetilde{\bm{\mu}}_t
    &=
    \widetilde{\beta}_t
    \left[
        \frac{\sqrt{\alpha_t}}{\beta_t}\bm{x}_t
        +
        \frac{
            \sqrt{\overline{\alpha}_{t-1}}
        }{
            1-\overline{\alpha}_{t-1}
        }\bm{x}_0
    \right]\\
    &=
    \frac{
        \sqrt{\alpha_t}
        (1-\overline{\alpha}_{t-1})
    }{
        1-\overline{\alpha}_t
    }\bm{x}_t
    +
    \frac{
        \sqrt{\overline{\alpha}_{t-1}}\beta_t
    }{
        1-\overline{\alpha}_t
    }\bm{x}_0.
\end{align*}
Это доказывает требуемую формулу.
\end{proof}

\subsection{Обратный генеративный процесс}

Генеративная модель задаётся как
\[
    p_\theta(\bm{x}_{0:T})
    =
    p(\bm{x}_T)
    \prod_{t=1}^{T}
    p_\theta(\bm{x}_{t-1}\mid\bm{x}_t),
\]
где
\[
    p(\bm{x}_T)=\N(\bm{0},\bm{I}),
\]
\[
    p_\theta(\bm{x}_{t-1}\mid\bm{x}_t)
    =
    \N\left(
        \bm{x}_{t-1};
        \bm{\mu}_\theta(\bm{x}_t,t),
        \bm{\Sigma}_\theta(\bm{x}_t,t)
    \right).
\]

Обратное распределение неизвестно, поскольку оно зависит от
распределения реальных данных. Его параметры аппроксимируются
нейронной сетью.

\subsection{Вариационная нижняя граница}

\begin{theorem}[Вариационная верхняя оценка отрицательного
логарифмического правдоподобия]
Для любого прямого распределения
$q(\bm{x}_{1:T}\mid\bm{x}_0)$ выполняется
\[
    -\log p_\theta(\bm{x}_0)
    \leq
    \E_q
    \left[
        \log
        \frac{
            q(\bm{x}_{1:T}\mid\bm{x}_0)
        }{
            p_\theta(\bm{x}_{0:T})
        }
    \right].
\]
Кроме того, правая часть раскладывается как
\begin{align*}
    \mathcal{L}_{\mathrm{VLB}}
    &=
    \KL\left(
        q(\bm{x}_T\mid\bm{x}_0)
        \,\|\,p(\bm{x}_T)
    \right)\\
    &\quad+
    \sum_{t=2}^{T}
    \E_{q(\bm{x}_t\mid\bm{x}_0)}
    \left[
        \KL\left(
            q(\bm{x}_{t-1}\mid\bm{x}_t,\bm{x}_0)
            \,\|\,
            p_\theta(\bm{x}_{t-1}\mid\bm{x}_t)
        \right)
    \right]\\
    &\quad-
    \E_{q(\bm{x}_1\mid\bm{x}_0)}
    \left[
        \log p_\theta(\bm{x}_0\mid\bm{x}_1)
    \right].
\end{align*}
\end{theorem}

\begin{proof}
Запишем маргинальное правдоподобие:
\[
    p_\theta(\bm{x}_0)
    =
    \int
    p_\theta(\bm{x}_{0:T})
    \dd\bm{x}_{1:T}.
\]
Умножим и разделим подынтегральное выражение на
$q(\bm{x}_{1:T}\mid\bm{x}_0)$:
\[
    p_\theta(\bm{x}_0)
    =
    \int
    q(\bm{x}_{1:T}\mid\bm{x}_0)
    \frac{
        p_\theta(\bm{x}_{0:T})
    }{
        q(\bm{x}_{1:T}\mid\bm{x}_0)
    }
    \dd\bm{x}_{1:T}.
\]
Следовательно,
\[
    p_\theta(\bm{x}_0)
    =
    \E_q
    \left[
        \frac{
            p_\theta(\bm{x}_{0:T})
        }{
            q(\bm{x}_{1:T}\mid\bm{x}_0)
        }
    \right].
\]
По неравенству Йенсена для вогнутой функции $\log$:
\[
    \log p_\theta(\bm{x}_0)
    \geq
    \E_q
    \left[
        \log
        \frac{
            p_\theta(\bm{x}_{0:T})
        }{
            q(\bm{x}_{1:T}\mid\bm{x}_0)
        }
    \right].
\]
После умножения на $-1$:
\[
    -\log p_\theta(\bm{x}_0)
    \leq
    \E_q
    \left[
        \log
        \frac{
            q(\bm{x}_{1:T}\mid\bm{x}_0)
        }{
            p_\theta(\bm{x}_{0:T})
        }
    \right].
\]

Теперь воспользуемся обратной факторизацией прямого процесса:
\[
    q(\bm{x}_{1:T}\mid\bm{x}_0)
    =
    q(\bm{x}_T\mid\bm{x}_0)
    \prod_{t=2}^{T}
    q(\bm{x}_{t-1}\mid\bm{x}_t,\bm{x}_0).
\]
Генеративный процесс факторизуется как
\[
    p_\theta(\bm{x}_{0:T})
    =
    p(\bm{x}_T)
    p_\theta(\bm{x}_0\mid\bm{x}_1)
    \prod_{t=2}^{T}
    p_\theta(\bm{x}_{t-1}\mid\bm{x}_t).
\]
Поэтому
\begin{align*}
    \log
    \frac{
        q(\bm{x}_{1:T}\mid\bm{x}_0)
    }{
        p_\theta(\bm{x}_{0:T})
    }
    &=
    \log
    \frac{
        q(\bm{x}_T\mid\bm{x}_0)
    }{
        p(\bm{x}_T)
    }\\
    &\quad+
    \sum_{t=2}^{T}
    \log
    \frac{
        q(\bm{x}_{t-1}\mid\bm{x}_t,\bm{x}_0)
    }{
        p_\theta(\bm{x}_{t-1}\mid\bm{x}_t)
    }\\
    &\quad-
    \log p_\theta(\bm{x}_0\mid\bm{x}_1).
\end{align*}
Взятие математического ожидания от первого отношения даёт
соответствующую KL-дивергенцию. Условное ожидание каждого слагаемого
суммы по $\bm{x}_{t-1}$ также даёт KL-дивергенцию. Последнее
слагаемое является ошибкой реконструкции. Получаем требуемое
разложение.
\end{proof}

Первое слагаемое не содержит обучаемых параметров, если расписание
шума фиксировано. Основное обучение связано с приближением точных
обратных переходов.

\subsection{Параметризации нейронной сети}

\subsubsection{Предсказание исходного объекта}

Сеть может предсказывать
\[
    \widehat{\bm{x}}_{0,\theta}
    =
    \bm{x}_\theta(\bm{x}_t,t).
\]

\subsubsection{Предсказание шума}

Из формулы зашумления
\[
    \bm{x}_t
    =
    \sqrt{\overline{\alpha}_t}\bm{x}_0
    +
    \sqrt{1-\overline{\alpha}_t}\bm{\eps}
\]
следует
\[
    \bm{x}_0
    =
    \frac{
        \bm{x}_t
        -
        \sqrt{1-\overline{\alpha}_t}\bm{\eps}
    }{
        \sqrt{\overline{\alpha}_t}
    }.
\]
Если сеть предсказывает шум
$\bm{\eps}_\theta(\bm{x}_t,t)$, то
\[
    \widehat{\bm{x}}_{0,\theta}
    =
    \frac{
        \bm{x}_t
        -
        \sqrt{1-\overline{\alpha}_t}
        \bm{\eps}_\theta(\bm{x}_t,t)
    }{
        \sqrt{\overline{\alpha}_t}
    }.
\]

Среднее обратного процесса удобно параметризовать как
\[
    \bm{\mu}_\theta(\bm{x}_t,t)
    =
    \frac{1}{\sqrt{\alpha_t}}
    \left(
        \bm{x}_t
        -
        \frac{\beta_t}
             {\sqrt{1-\overline{\alpha}_t}}
        \bm{\eps}_\theta(\bm{x}_t,t)
    \right).
\]

\subsubsection{\(v\)-параметризация}

Положим
\[
    a_t=\sqrt{\overline{\alpha}_t},
    \qquad
    b_t=\sqrt{1-\overline{\alpha}_t}.
\]
Тогда
\[
    \bm{x}_t=a_t\bm{x}_0+b_t\bm{\eps}.
\]
Определим
\[
    \bm{v}_t=a_t\bm{\eps}-b_t\bm{x}_0.
\]
Поскольку $a_t^2+b_t^2=1$, обратное преобразование имеет вид
\[
    \bm{x}_0=a_t\bm{x}_t-b_t\bm{v}_t,
\]
\[
    \bm{\eps}=b_t\bm{x}_t+a_t\bm{v}_t.
\]
Сеть может обучаться предсказывать $\bm{v}_t$ вместо
$\bm{x}_0$ или $\bm{\eps}$.

\subsection{Связь вариационной цели с MSE шума}

\begin{theorem}[Квадратичная ошибка предсказания шума]
Пусть
\[
    p_\theta(\bm{x}_{t-1}\mid\bm{x}_t)
    =
    \N(
        \bm{\mu}_\theta(\bm{x}_t,t),
        \sigma_t^2\bm{I}
    ),
\]
где $\sigma_t^2$ не зависит от $\theta$, и
\[
    \bm{\mu}_\theta(\bm{x}_t,t)
    =
    \frac{1}{\sqrt{\alpha_t}}
    \left(
        \bm{x}_t
        -
        \frac{\beta_t}
             {\sqrt{1-\overline{\alpha}_t}}
        \bm{\eps}_\theta(\bm{x}_t,t)
    \right).
\]
Тогда зависящая от $\theta$ часть
\[
    \KL\left(
        q(\bm{x}_{t-1}\mid\bm{x}_t,\bm{x}_0)
        \,\|\,
        p_\theta(\bm{x}_{t-1}\mid\bm{x}_t)
    \right)
\]
равна
\[
    \frac{
        \beta_t^2
    }{
        2\sigma_t^2\alpha_t
        (1-\overline{\alpha}_t)
    }
    \left\|
        \bm{\eps}
        -
        \bm{\eps}_\theta(\bm{x}_t,t)
    \right\|^2
\]
с точностью до слагаемого, не зависящего от $\theta$.
\end{theorem}

\begin{proof}
Точное среднее апостериорного перехода можно переписать через шум.
Так как
\[
    \bm{x}_t
    =
    \sqrt{\overline{\alpha}_t}\bm{x}_0
    +
    \sqrt{1-\overline{\alpha}_t}\bm{\eps},
\]
имеем
\[
    \bm{x}_0
    =
    \frac{
        \bm{x}_t
        -
        \sqrt{1-\overline{\alpha}_t}\bm{\eps}
    }{
        \sqrt{\overline{\alpha}_t}
    }.
\]
Подстановка в формулу для
$\widetilde{\bm{\mu}}_t(\bm{x}_t,\bm{x}_0)$ и упрощение дают
\[
    \widetilde{\bm{\mu}}_t
    =
    \frac{1}{\sqrt{\alpha_t}}
    \left(
        \bm{x}_t
        -
        \frac{\beta_t}
             {\sqrt{1-\overline{\alpha}_t}}
        \bm{\eps}
    \right).
\]

Для двух нормальных распределений
\[
    q=\N(\bm{\mu}_q,\bm{\Sigma}_q),
    \qquad
    p=\N(\bm{\mu}_p,\bm{\Sigma}_p)
\]
KL-дивергенция содержит зависящее от средних слагаемое
\[
    \frac{1}{2}
    (\bm{\mu}_p-\bm{\mu}_q)^\top
    \bm{\Sigma}_p^{-1}
    (\bm{\mu}_p-\bm{\mu}_q).
\]
В рассматриваемом случае
\[
    \bm{\Sigma}_p=\sigma_t^2\bm{I},
    \qquad
    \bm{\Sigma}_p^{-1}
    =
    \frac{1}{\sigma_t^2}\bm{I}.
\]
Поэтому зависящая от $\theta$ часть равна
\[
    \frac{1}{2\sigma_t^2}
    \|\bm{\mu}_\theta-\widetilde{\bm{\mu}}_t\|^2.
\]

Разность средних:
\begin{align*}
    \bm{\mu}_\theta-\widetilde{\bm{\mu}}_t
    &=
    \frac{1}{\sqrt{\alpha_t}}
    \left(
        -
        \frac{\beta_t}
             {\sqrt{1-\overline{\alpha}_t}}
        \bm{\eps}_\theta
        +
        \frac{\beta_t}
             {\sqrt{1-\overline{\alpha}_t}}
        \bm{\eps}
    \right)\\
    &=
    \frac{\beta_t}
         {\sqrt{\alpha_t}
          \sqrt{1-\overline{\alpha}_t}}
    (\bm{\eps}-\bm{\eps}_\theta).
\end{align*}
Следовательно,
\[
    \frac{1}{2\sigma_t^2}
    \|\bm{\mu}_\theta-\widetilde{\bm{\mu}}_t\|^2
    =
    \frac{
        \beta_t^2
    }{
        2\sigma_t^2\alpha_t
        (1-\overline{\alpha}_t)
    }
    \|\bm{\eps}-\bm{\eps}_\theta\|^2.
\]
Оставшиеся члены KL-дивергенции зависят от ковариаций, но не от
$\theta$ при фиксированной $\sigma_t^2$.
\end{proof}

На практике часто используется упрощённая невзвешенная функция потерь
\[
    \mathcal{L}_{\text{simple}}(\theta)
    =
    \E_{
        \bm{x}_0,\,
        t,\,
        \bm{\eps}
    }
    \left[
        \left\|
            \bm{\eps}
            -
            \bm{\eps}_\theta(\bm{x}_t,t)
        \right\|_2^2
    \right],
\]
где
\[
    \bm{x}_t
    =
    \sqrt{\overline{\alpha}_t}\bm{x}_0
    +
    \sqrt{1-\overline{\alpha}_t}\bm{\eps}.
\]

\subsection{Алгоритм обучения DDPM}

\begin{algorithmbox}{Обучение модели предсказания шума}
    \item Выбрать объект
    \[
        \bm{x}_0\sim q_{\text{data}}.
    \]
    \item Выбрать момент времени, например
    \[
        t\sim\operatorname{Uniform}\{1,\ldots,T\}.
    \]
    \item Сгенерировать шум
    \[
        \bm{\eps}\sim\N(\bm{0},\bm{I}).
    \]
    \item Получить зашумлённый объект напрямую:
    \[
        \bm{x}_t
        =
        \sqrt{\overline{\alpha}_t}\bm{x}_0
        +
        \sqrt{1-\overline{\alpha}_t}\bm{\eps}.
    \]
    \item Передать $\bm{x}_t$ и эмбеддинг времени $t$ в нейронную сеть:
    \[
        \widehat{\bm{\eps}}
        =
        \bm{\eps}_\theta(\bm{x}_t,t).
    \]
    \item Вычислить ошибку
    \[
        \mathcal{L}
        =
        \|\bm{\eps}-\widehat{\bm{\eps}}\|^2.
    \]
    \item Обновить параметры $\theta$ методом стохастического
    градиентного спуска.
\end{algorithmbox}

Для одного обучающего примера вычисляется только один случайный уровень
шума. Поэтому стоимость шага обучения не пропорциональна $T$.

\subsection{Алгоритм генерации DDPM}

\begin{algorithmbox}{Стохастическая генерация DDPM}
    \item Сгенерировать начальный шум
    \[
        \bm{x}_T\sim\N(\bm{0},\bm{I}).
    \]
    \item Для $t=T,T-1,\ldots,1$ выполнить следующие шаги.
    \item Предсказать шум:
    \[
        \widehat{\bm{\eps}}
        =
        \bm{\eps}_\theta(\bm{x}_t,t).
    \]
    \item Вычислить среднее обратного перехода:
    \[
        \bm{\mu}_\theta
        =
        \frac{1}{\sqrt{\alpha_t}}
        \left(
            \bm{x}_t
            -
            \frac{\beta_t}
                 {\sqrt{1-\overline{\alpha}_t}}
            \widehat{\bm{\eps}}
        \right).
    \]
    \item Если $t>1$, сгенерировать
    \[
        \bm{z}\sim\N(\bm{0},\bm{I}),
    \]
    иначе положить $\bm{z}=\bm{0}$.
    \item Выполнить переход
    \[
        \bm{x}_{t-1}
        =
        \bm{\mu}_\theta
        +
        \sigma_t\bm{z}.
    \]
    Частый выбор:
    \[
        \sigma_t^2=\widetilde{\beta}_t.
    \]
    \item Вернуть итоговый объект $\bm{x}_0$.
\end{algorithmbox}

Обучение допускает параллельную обработку объектов мини-пакета, но
стандартная генерация последовательна по времени: следующий шаг
зависит от результата предыдущего.

\subsection{Минимальный одномерный пример}

Пусть
\[
    x_0=2,\qquad
    \beta_1=0.1,\qquad
    \beta_2=0.2.
\]
Тогда
\[
    \alpha_1=0.9,\qquad
    \alpha_2=0.8,
\]
\[
    \overline{\alpha}_1=0.9,\qquad
    \overline{\alpha}_2=0.72.
\]

Возьмём шум $\eps=-0.5$. Прямое зашумление до момента $t=2$:
\begin{align*}
    x_2
    &=
    \sqrt{0.72}\cdot 2
    +
    \sqrt{0.28}\cdot(-0.5)\\
    &\approx
    1.6971-0.2646\\
    &\approx 1.4325.
\end{align*}

Пусть нейронная сеть предсказала
\[
    \eps_\theta(x_2,2)=-0.45.
\]
Среднее обратного перехода:
\begin{align*}
    \mu_\theta
    &=
    \frac{1}{\sqrt{0.8}}
    \left(
        1.4325
        -
        \frac{0.2}{\sqrt{0.28}}(-0.45)
    \right)\\
    &\approx
    1.7918.
\end{align*}
Дисперсия точного апостериорного перехода:
\[
    \widetilde{\beta}_2
    =
    \frac{1-0.9}{1-0.72}\cdot 0.2
    \approx 0.07143.
\]
Следовательно,
\[
    x_1
    =
    1.7918+\sqrt{0.07143}\,z,
    \qquad
    z\sim\N(0,1).
\]
Если рассматривать только среднее и взять $z=0$, получится
$x_1\approx1.7918$.

\subsection{Score-функция и denoising score matching}

\begin{definition}
Score-функцией плотности $p(\bm{x})$ называется градиент логарифма
плотности по данным:
\[
    \bm{s}_p(\bm{x})
    =
    \nabla_{\bm{x}}\log p(\bm{x}).
\]
\end{definition}

Для условного нормального распределения
\[
    q(\bm{x}_t\mid\bm{x}_0)
    =
    \N\left(
        \sqrt{\overline{\alpha}_t}\bm{x}_0,
        (1-\overline{\alpha}_t)\bm{I}
    \right)
\]
условная score-функция равна
\[
    \nabla_{\bm{x}_t}
    \log q(\bm{x}_t\mid\bm{x}_0)
    =
    -
    \frac{
        \bm{x}_t-\sqrt{\overline{\alpha}_t}\bm{x}_0
    }{
        1-\overline{\alpha}_t
    }.
\]
Поскольку
\[
    \bm{x}_t-\sqrt{\overline{\alpha}_t}\bm{x}_0
    =
    \sqrt{1-\overline{\alpha}_t}\bm{\eps},
\]
получаем
\[
    \nabla_{\bm{x}_t}
    \log q(\bm{x}_t\mid\bm{x}_0)
    =
    -
    \frac{\bm{\eps}}
         {\sqrt{1-\overline{\alpha}_t}}.
\]
Следовательно, модель шума задаёт модель score-функции:
\[
    \bm{s}_\theta(\bm{x}_t,t)
    =
    -
    \frac{
        \bm{\eps}_\theta(\bm{x}_t,t)
    }{
        \sqrt{1-\overline{\alpha}_t}
    }.
\]

\begin{theorem}[Эквивалентность условного и маргинального
denoising score matching]
Пусть случайная величина $\bm{X}$ имеет плотность $p(\bm{x})$, а
$\bm{Y}$ при условии $\bm{X}=\bm{x}$ имеет дифференцируемую плотность
$q(\bm{y}\mid\bm{x})$. Пусть
\[
    p_Y(\bm{y})
    =
    \int p(\bm{x})q(\bm{y}\mid\bm{x})\dd\bm{x}
\]
--- маргинальная плотность зашумлённых данных. При условиях,
позволяющих менять местами дифференцирование и интегрирование,
для любой измеримой векторной функции $\bm{s}(\bm{y})$ с конечными
вторыми моментами выполняется
\begin{align*}
    &\E
    \left[
        \left\|
            \bm{s}(\bm{Y})
            -
            \nabla_{\bm{y}}
            \log q(\bm{Y}\mid\bm{X})
        \right\|^2
    \right]\\
    &\qquad=
    \E
    \left[
        \left\|
            \bm{s}(\bm{Y})
            -
            \nabla_{\bm{y}}\log p_Y(\bm{Y})
        \right\|^2
    \right]
    +C,
\end{align*}
где $C$ не зависит от $\bm{s}$. Следовательно, обе функции потерь
имеют один и тот же оптимизатор:
\[
    \bm{s}^{*}(\bm{y})
    =
    \nabla_{\bm{y}}\log p_Y(\bm{y}).
\]
\end{theorem}

\begin{proof}
Сначала вычислим условное ожидание условной score-функции. Имеем
\[
    \nabla_{\bm{y}}p_Y(\bm{y})
    =
    \nabla_{\bm{y}}
    \int
    p(\bm{x})q(\bm{y}\mid\bm{x})
    \dd\bm{x}.
\]
При сделанных предположениях:
\[
    \nabla_{\bm{y}}p_Y(\bm{y})
    =
    \int
    p(\bm{x})
    \nabla_{\bm{y}}q(\bm{y}\mid\bm{x})
    \dd\bm{x}.
\]
Используя тождество
\[
    \nabla_{\bm{y}}q(\bm{y}\mid\bm{x})
    =
    q(\bm{y}\mid\bm{x})
    \nabla_{\bm{y}}\log q(\bm{y}\mid\bm{x}),
\]
получаем
\begin{align*}
    \nabla_{\bm{y}}p_Y(\bm{y})
    &=
    \int
    p(\bm{x})q(\bm{y}\mid\bm{x})
    \nabla_{\bm{y}}\log q(\bm{y}\mid\bm{x})
    \dd\bm{x}\\
    &=
    p_Y(\bm{y})
    \E\left[
        \nabla_{\bm{y}}\log q(\bm{Y}\mid\bm{X})
        \given
        \bm{Y}=\bm{y}
    \right].
\end{align*}
После деления на $p_Y(\bm{y})$:
\[
    \E\left[
        \nabla_{\bm{y}}\log q(\bm{Y}\mid\bm{X})
        \given
        \bm{Y}=\bm{y}
    \right]
    =
    \nabla_{\bm{y}}\log p_Y(\bm{y}).
\]

Обозначим
\[
    \bm{Z}
    =
    \nabla_{\bm{y}}\log q(\bm{Y}\mid\bm{X}),
\]
\[
    \bm{m}(\bm{Y})
    =
    \E[\bm{Z}\mid\bm{Y}]
    =
    \nabla_{\bm{y}}\log p_Y(\bm{Y}).
\]
Разложим ошибку:
\[
    \bm{s}(\bm{Y})-\bm{Z}
    =
    \bigl(\bm{s}(\bm{Y})-\bm{m}(\bm{Y})\bigr)
    +
    \bigl(\bm{m}(\bm{Y})-\bm{Z}\bigr).
\]
Квадрат нормы равен
\begin{align*}
    \|\bm{s}-\bm{Z}\|^2
    &=
    \|\bm{s}-\bm{m}\|^2
    +
    \|\bm{m}-\bm{Z}\|^2\\
    &\quad+
    2(\bm{s}-\bm{m})^\top(\bm{m}-\bm{Z}).
\end{align*}
Условное ожидание перекрёстного слагаемого при фиксированном $\bm{Y}$
равно нулю:
\begin{align*}
    &\E\left[
        (\bm{s}-\bm{m})^\top(\bm{m}-\bm{Z})
        \given \bm{Y}
    \right]\\
    &\quad=
    (\bm{s}(\bm{Y})-\bm{m}(\bm{Y}))^\top
    \left(
        \bm{m}(\bm{Y})
        -
        \E[\bm{Z}\mid\bm{Y}]
    \right)
    =0.
\end{align*}
Поэтому
\begin{align*}
    \E\|\bm{s}(\bm{Y})-\bm{Z}\|^2
    &=
    \E\|\bm{s}(\bm{Y})-\bm{m}(\bm{Y})\|^2\\
    &\quad+
    \E\|\bm{m}(\bm{Y})-\bm{Z}\|^2.
\end{align*}
Второе слагаемое не зависит от $\bm{s}$ и является константой $C$.
Подставляя
\[
    \bm{m}(\bm{Y})
    =
    \nabla_{\bm{y}}\log p_Y(\bm{Y}),
\]
получаем утверждение теоремы.
\end{proof}

Таким образом, предсказание искусственно добавленного шума обучает
градиент логарифма маргинальной плотности зашумлённых данных.

\subsection{Непрерывное время и стохастические дифференциальные
уравнения}

Дискретную диффузию можно обобщить на стохастическое
дифференциальное уравнение:
\[
    \dd\bm{X}_t
    =
    \bm{f}(\bm{X}_t,t)\dd t
    +
    g(t)\dd\bm{W}_t,
\]
где $\bm{W}_t$ --- винеровский процесс.

Плотность $p_t(\bm{x})$ удовлетворяет уравнению Фоккера--Планка:
\[
    \frac{\partial p_t}{\partial t}
    =
    -
    \nabla\cdot(\bm{f}p_t)
    +
    \frac{1}{2}g(t)^2\Delta p_t.
\]

\begin{theorem}[Обратное по времени SDE]
Пусть процесс $\bm{X}_t$, $t\in[0,T]$, удовлетворяет
\[
    \dd\bm{X}_t
    =
    \bm{f}(\bm{X}_t,t)\dd t
    +
    g(t)\dd\bm{W}_t,
\]
где $g(t)>0$ не зависит от $\bm{x}$. Предположим, что плотности
$p_t(\bm{x})$ положительны, достаточно гладкие, а соответствующие
стохастические уравнения имеют единственное распределение решения.

Для обратного времени $s=T-t$ процесс
\[
    \bm{Y}_s=\bm{X}_{T-s}
\]
имеет SDE
\[
    \dd\bm{Y}_s
    =
    \left[
        -\bm{f}(\bm{Y}_s,T-s)
        +
        g(T-s)^2
        \nabla_{\bm{y}}\log p_{T-s}(\bm{Y}_s)
    \right]\dd s
    +
    g(T-s)\dd\overline{\bm{W}}_s.
\]
Эквивалентно, если записывать уравнение в переменной $t$, движущейся
от $T$ к $0$:
\[
    \dd\bm{X}_t
    =
    \left[
        \bm{f}(\bm{X}_t,t)
        -
        g(t)^2\nabla_{\bm{x}}\log p_t(\bm{X}_t)
    \right]\dd t
    +
    g(t)\dd\overline{\bm{W}}_t,
\]
где интегрирование выполняется назад по времени.
\end{theorem}

\begin{proof}
Обозначим плотность обратного процесса через
\[
    q_s(\bm{y})=p_{T-s}(\bm{y}).
\]
Тогда
\[
    \frac{\partial q_s}{\partial s}
    =
    -
    \left.
    \frac{\partial p_t}{\partial t}
    \right|_{t=T-s}.
\]
Из уравнения Фоккера--Планка:
\[
    \frac{\partial q_s}{\partial s}
    =
    \nabla\cdot
    \bigl(
        \bm{f}(\bm{y},T-s)q_s(\bm{y})
    \bigr)
    -
    \frac{1}{2}g(T-s)^2\Delta q_s(\bm{y}).
\]

Рассмотрим SDE с обратным дрейфом
\[
    \bm{b}_{\mathrm{rev}}(\bm{y},s)
    =
    -\bm{f}(\bm{y},T-s)
    +
    g(T-s)^2\nabla_{\bm{y}}\log q_s(\bm{y}).
\]
Его уравнение Фоккера--Планка:
\[
    \frac{\partial q_s}{\partial s}
    =
    -
    \nabla\cdot(\bm{b}_{\mathrm{rev}}q_s)
    +
    \frac{1}{2}g(T-s)^2\Delta q_s.
\]
Подставим дрейф:
\begin{align*}
    -
    \nabla\cdot(\bm{b}_{\mathrm{rev}}q_s)
    &=
    -
    \nabla\cdot
    \left(
        -\bm{f}q_s
        +
        g^2q_s\nabla\log q_s
    \right)\\
    &=
    \nabla\cdot(\bm{f}q_s)
    -
    g^2\nabla\cdot(q_s\nabla\log q_s).
\end{align*}
Но
\[
    q_s\nabla\log q_s=\nabla q_s,
\]
поэтому
\[
    \nabla\cdot(q_s\nabla\log q_s)=\Delta q_s.
\]
Следовательно,
\begin{align*}
    \frac{\partial q_s}{\partial s}
    &=
    \nabla\cdot(\bm{f}q_s)
    -
    g^2\Delta q_s
    +
    \frac{1}{2}g^2\Delta q_s\\
    &=
    \nabla\cdot(\bm{f}q_s)
    -
    \frac{1}{2}g^2\Delta q_s.
\end{align*}
Это в точности уравнение эволюции плотности
$p_{T-s}$. При предположении единственности распределения решения
указанный SDE задаёт обратный по времени процесс.
\end{proof}

Для генерации необходимо аппроксимировать неизвестную score-функцию:
\[
    \nabla_{\bm{x}}\log p_t(\bm{x})
    \approx
    \bm{s}_\theta(\bm{x},t).
\]

\subsection{Probability flow ODE}

\begin{theorem}[Детерминированное уравнение с теми же маргиналами]
В условиях предыдущей теоремы обыкновенное дифференциальное уравнение
\[
    \frac{\dd\bm{X}_t}{\dd t}
    =
    \bm{f}(\bm{X}_t,t)
    -
    \frac{1}{2}g(t)^2
    \nabla_{\bm{x}}\log p_t(\bm{X}_t)
\]
имеет те же маргинальные плотности $p_t$, что и прямое SDE
\[
    \dd\bm{X}_t
    =
    \bm{f}(\bm{X}_t,t)\dd t
    +
    g(t)\dd\bm{W}_t.
\]
\end{theorem}

\begin{proof}
Плотность детерминированного потока с полем скоростей
\[
    \bm{v}(\bm{x},t)
    =
    \bm{f}(\bm{x},t)
    -
    \frac{1}{2}g(t)^2\nabla\log p_t(\bm{x})
\]
удовлетворяет уравнению непрерывности
\[
    \frac{\partial p_t}{\partial t}
    =
    -\nabla\cdot(\bm{v}p_t).
\]
Подставим $\bm{v}$:
\begin{align*}
    -\nabla\cdot(\bm{v}p_t)
    &=
    -\nabla\cdot(\bm{f}p_t)
    +
    \frac{1}{2}g(t)^2
    \nabla\cdot
    \left(
        p_t\nabla\log p_t
    \right)\\
    &=
    -\nabla\cdot(\bm{f}p_t)
    +
    \frac{1}{2}g(t)^2
    \nabla\cdot(\nabla p_t)\\
    &=
    -\nabla\cdot(\bm{f}p_t)
    +
    \frac{1}{2}g(t)^2\Delta p_t.
\end{align*}
Получено то же уравнение Фоккера--Планка, что и для исходного SDE.
При одинаковом начальном распределении и единственности решения
плотности совпадают в каждый момент времени.
\end{proof}

Probability flow ODE позволяет выполнять детерминированную генерацию и
связывает диффузионные модели с непрерывными нормализующими потоками.

\subsection{DDPM как дискретизация VP-SDE}

Рассмотрим variance-preserving SDE:
\[
    \dd\bm{X}_t
    =
    -\frac{1}{2}\beta(t)\bm{X}_t\dd t
    +
    \sqrt{\beta(t)}\dd\bm{W}_t.
\]
Его условное распределение имеет вид
\[
    \bm{X}_t
    =
    \exp\left(
        -\frac{1}{2}\int_0^t\beta(s)\dd s
    \right)\bm{X}_0
    +
    \sqrt{
        1-
        \exp\left(
            -\int_0^t\beta(s)\dd s
        \right)
    }\bm{\eps}.
\]
Это непрерывный аналог формулы
\[
    \bm{x}_t
    =
    \sqrt{\overline{\alpha}_t}\bm{x}_0
    +
    \sqrt{1-\overline{\alpha}_t}\bm{\eps}.
\]

\subsection{DDIM}

DDPM использует стохастические обратные переходы. DDIM задаёт семейство
процессов, совместимых с той же обученной моделью, но допускающих
детерминированную генерацию.

Пусть
\[
    T=\tau_N>\tau_{N-1}>\cdots>\tau_1>\tau_0=0
\]
--- выбранная подпоследовательность временных шагов. Для перехода
$t=\tau_i$ в $s=\tau_{i-1}$ вычисляется
\[
    \widehat{\bm{x}}_0
    =
    \frac{
        \bm{x}_t
        -
        \sqrt{1-\overline{\alpha}_t}
        \bm{\eps}_\theta(\bm{x}_t,t)
    }{
        \sqrt{\overline{\alpha}_t}
    }.
\]
Параметр стохастичности:
\[
    \sigma_{t\rightarrow s}
    =
    \eta
    \sqrt{
        \frac{1-\overline{\alpha}_s}
             {1-\overline{\alpha}_t}
        \left(
            1-\frac{\overline{\alpha}_t}
                    {\overline{\alpha}_s}
        \right)
    }.
\]
Переход:
\begin{align*}
    \bm{x}_s
    &=
    \sqrt{\overline{\alpha}_s}
    \widehat{\bm{x}}_0\\
    &\quad+
    \sqrt{
        1-\overline{\alpha}_s
        -
        \sigma_{t\rightarrow s}^2
    }
    \bm{\eps}_\theta(\bm{x}_t,t)
    +
    \sigma_{t\rightarrow s}\bm{z}.
\end{align*}

При $\eta=0$ процесс детерминирован:
\[
    \bm{x}_s
    =
    \sqrt{\overline{\alpha}_s}
    \widehat{\bm{x}}_0
    +
    \sqrt{1-\overline{\alpha}_s}
    \bm{\eps}_\theta(\bm{x}_t,t).
\]

\begin{algorithmbox}{Ускоренная генерация DDIM}
    \item Выбрать небольшое число временных точек
    \[
        T=\tau_N>\cdots>\tau_0=0.
    \]
    \item Сгенерировать
    \[
        \bm{x}_{\tau_N}\sim\N(\bm{0},\bm{I}).
    \]
    \item Для $i=N,\ldots,1$ положить
    $t=\tau_i$, $s=\tau_{i-1}$.
    \item Предсказать шум
    \[
        \widehat{\bm{\eps}}
        =
        \bm{\eps}_\theta(\bm{x}_t,t).
    \]
    \item Оценить чистый объект
    \[
        \widehat{\bm{x}}_0
        =
        \frac{
            \bm{x}_t
            -
            \sqrt{1-\overline{\alpha}_t}\widehat{\bm{\eps}}
        }{
            \sqrt{\overline{\alpha}_t}
        }.
    \]
    \item Вычислить $\sigma_{t\rightarrow s}$ и получить $\bm{x}_s$
    по формуле DDIM.
    \item Вернуть $\bm{x}_0$.
\end{algorithmbox}

Уменьшение числа временных точек ускоряет генерацию, но слишком
крупный шаг может увеличить ошибку численного решения.

\subsection{Архитектура денойзера}

\subsubsection{U-Net}

Для изображений часто применяется U-Net:
\begin{enumerate}
    \item сжимающая часть уменьшает пространственное разрешение и
    увеличивает число каналов;
    \item центральные блоки обрабатывают признаки низкого разрешения;
    \item расширяющая часть восстанавливает разрешение;
    \item skip connections передают детальные признаки между
    симметричными уровнями.
\end{enumerate}

Типичный диффузионный U-Net содержит:
\begin{itemize}
    \item остаточные свёрточные блоки;
    \item нормализацию;
    \item эмбеддинг времени;
    \item self-attention на части разрешений;
    \item cross-attention для условной генерации.
\end{itemize}

\subsubsection{Эмбеддинг времени}

Момент $t$ преобразуется в вектор, например синусоидальным
кодированием:
\[
    e_{2i}(t)
    =
    \sin(\omega_i t),
    \qquad
    e_{2i+1}(t)
    =
    \cos(\omega_i t).
\]
После небольшой MLP этот вектор добавляется или подмешивается в
промежуточные блоки сети. Без времени одна и та же сеть не могла бы
различать слабый и сильный уровни шума.

\subsubsection{Diffusion Transformer}

Вместо U-Net можно использовать трансформер:
\begin{enumerate}
    \item изображение или латентное представление разбивается на
    патчи;
    \item каждый патч преобразуется в токен;
    \item к токенам добавляются позиционные признаки;
    \item время и условие передаются через дополнительные токены,
    адаптивную нормализацию или cross-attention;
    \item трансформер предсказывает шум, score-функцию или
    $\bm{v}$-параметризацию для всех патчей.
\end{enumerate}

Таким образом, понятия ``трансформер'' и ``диффузионная модель'' не
конкурируют: трансформер может быть внутренней архитектурой
диффузионного денойзера.

\subsection{Условная генерация}

Пусть $\bm{c}$ --- условие: класс, текст, изображение, маска,
аудиосигнал или другой объект. Тогда сеть имеет вид
\[
    \bm{\eps}_\theta(\bm{x}_t,t,\bm{c}).
\]

Для текстового условия часто используется отдельный трансформерный
кодировщик:
\[
    \bm{C}
    =
    \operatorname{TextEncoder}(\bm{c}).
\]
В денойзере выполняется cross-attention:
\[
    \bm{Q}=\bm{X}_{\text{image}}\bm{W}^Q,
    \qquad
    \bm{K}=\bm{C}\bm{W}^K,
    \qquad
    \bm{V}=\bm{C}\bm{W}^V.
\]

\subsection{Classifier guidance}

По формуле Байеса:
\[
    p_t(\bm{x}\mid y)
    \propto
    p_t(\bm{x})p_t(y\mid\bm{x}).
\]
Следовательно,
\[
    \nabla_{\bm{x}}\log p_t(\bm{x}\mid y)
    =
    \nabla_{\bm{x}}\log p_t(\bm{x})
    +
    \nabla_{\bm{x}}\log p_t(y\mid\bm{x}).
\]
Условную score-функцию можно приближать как
\[
    \bm{s}_{\text{guided}}
    =
    \bm{s}_\theta(\bm{x}_t,t)
    +
    w\nabla_{\bm{x}_t}
    \log p_\phi(y\mid\bm{x}_t,t),
\]
где $p_\phi$ --- отдельно обученный классификатор зашумлённых данных,
а $w$ --- сила управления.

Недостатки:
\begin{itemize}
    \item требуется дополнительный классификатор;
    \item классификатор должен работать на всех уровнях шума;
    \item метод не всегда удобен для сложных условий, например текста.
\end{itemize}

\subsection{Classifier-free guidance}

Одна модель обучается и в условном, и в безусловном режимах. Во время
обучения условие с некоторой вероятностью заменяется пустым:
\[
    \bm{c}
    \leftarrow
    \varnothing.
\]
На этапе генерации выполняются два предсказания:
\[
    \bm{\eps}_{\text{cond}}
    =
    \bm{\eps}_\theta(\bm{x}_t,t,\bm{c}),
\]
\[
    \bm{\eps}_{\text{uncond}}
    =
    \bm{\eps}_\theta(\bm{x}_t,t,\varnothing).
\]
Они объединяются:
\[
    \bm{\eps}_{\text{CFG}}
    =
    \bm{\eps}_{\text{uncond}}
    +
    w
    \left(
        \bm{\eps}_{\text{cond}}
        -
        \bm{\eps}_{\text{uncond}}
    \right).
\]

Интерпретация параметра:
\[
    w=0
    \quad\Rightarrow\quad
    \text{безусловная генерация},
\]
\[
    w=1
    \quad\Rightarrow\quad
    \text{обычное условное предсказание},
\]
\[
    w>1
    \quad\Rightarrow\quad
    \text{усиленное следование условию}.
\]

Слишком большое $w$ может снижать разнообразие, усиливать артефакты и
выводить модель за область распределения, встречавшуюся при обучении.

\begin{algorithmbox}{Шаг генерации с classifier-free guidance}
    \item Вычислить безусловное предсказание:
    \[
        \bm{\eps}_u
        =
        \bm{\eps}_\theta(\bm{x}_t,t,\varnothing).
    \]
    \item Вычислить условное предсказание:
    \[
        \bm{\eps}_c
        =
        \bm{\eps}_\theta(\bm{x}_t,t,\bm{c}).
    \]
    \item Объединить предсказания:
    \[
        \widehat{\bm{\eps}}
        =
        \bm{\eps}_u+w(\bm{\eps}_c-\bm{\eps}_u).
    \]
    \item Подставить $\widehat{\bm{\eps}}$ в формулу DDPM, DDIM или
    другого численного решателя.
\end{algorithmbox}

\subsection{Латентные диффузионные модели}

Пиксельная диффузия выполняет все шаги в пространстве изображений
высокого разрешения, что вычислительно дорого.

В латентной диффузии сначала обучается автоэнкодер:
\[
    \bm{z}_0=E(\bm{x}_0),
    \qquad
    \widehat{\bm{x}}_0=D(\bm{z}_0),
\]
где $E$ --- кодировщик, а $D$ --- декодировщик.

Диффузионный процесс обучается в латентном пространстве:
\[
    \bm{z}_t
    =
    \sqrt{\overline{\alpha}_t}\bm{z}_0
    +
    \sqrt{1-\overline{\alpha}_t}\bm{\eps}.
\]
После обратной диффузии итоговый латентный код декодируется:
\[
    \widehat{\bm{x}}=D(\widehat{\bm{z}}_0).
\]

Преимущества:
\begin{itemize}
    \item меньшее пространственное разрешение;
    \item более низкие затраты памяти;
    \item более быстрое обучение и сэмплирование;
    \item удаление части малозначимой пиксельной избыточности.
\end{itemize}

Недостаток заключается в том, что качество ограничивается
автоэнкодером: потерянные при кодировании детали диффузионная модель
обычно не может точно восстановить.

\subsection{Расписания шума}

Расписание определяет значения $\beta_t$ или, эквивалентно,
отношение сигнал/шум.

Для
\[
    \bm{x}_t
    =
    \sqrt{\overline{\alpha}_t}\bm{x}_0
    +
    \sqrt{1-\overline{\alpha}_t}\bm{\eps}
\]
отношение сигнал/шум можно определить как
\[
    \operatorname{SNR}(t)
    =
    \frac{\overline{\alpha}_t}
         {1-\overline{\alpha}_t}.
\]

Используются:
\begin{itemize}
    \item линейные расписания $\beta_t$;
    \item косинусные расписания $\overline{\alpha}_t$;
    \item непрерывные расписания по $\sigma$ или log-SNR;
    \item адаптированные веса функции потерь для разных уровней шума.
\end{itemize}

Слишком быстрое добавление шума может уничтожить полезный сигнал за
малое число шагов, а слишком медленное увеличивает число необходимых
итераций.

\subsection{Сравнение основных способов предсказания}

\begin{center}
\begin{tabular}{p{0.16\textwidth}p{0.28\textwidth}p{0.46\textwidth}}
\toprule
\textbf{Цель сети} &
\textbf{Определение} &
\textbf{Восстановление других величин}\\
\midrule
$\bm{x}_0$ &
Сеть предсказывает чистый объект &
\[
    \widehat{\bm{\eps}}
    =
    \frac{
        \bm{x}_t-\sqrt{\overline{\alpha}_t}\widehat{\bm{x}}_0
    }{
        \sqrt{1-\overline{\alpha}_t}
    }.
\]\\
\addlinespace
$\bm{\eps}$ &
Сеть предсказывает добавленный шум &
\[
    \widehat{\bm{x}}_0
    =
    \frac{
        \bm{x}_t
        -
        \sqrt{1-\overline{\alpha}_t}\widehat{\bm{\eps}}
    }{
        \sqrt{\overline{\alpha}_t}
    }.
\]\\
\addlinespace
$\bm{v}$ &
\[
    \bm{v}
    =
    \sqrt{\overline{\alpha}_t}\bm{\eps}
    -
    \sqrt{1-\overline{\alpha}_t}\bm{x}_0
\]
&
\[
    \widehat{\bm{x}}_0
    =
    \sqrt{\overline{\alpha}_t}\bm{x}_t
    -
    \sqrt{1-\overline{\alpha}_t}\widehat{\bm{v}},
\]
\[
    \widehat{\bm{\eps}}
    =
    \sqrt{1-\overline{\alpha}_t}\bm{x}_t
    +
    \sqrt{\overline{\alpha}_t}\widehat{\bm{v}}.
\]\\
\bottomrule
\end{tabular}
\end{center}

\subsection{Достоинства и ограничения диффузионных моделей}

\paragraph{Достоинства.}
\begin{itemize}
    \item Стабильная регрессионная функция потерь.
    \item Отсутствие обязательной состязательной оптимизации.
    \item Хорошее покрытие распределения данных.
    \item Гибкая условная генерация.
    \item Возможность редактирования, восстановления пропусков,
    суперразрешения и решения обратных задач.
    \item Связь с вариационным выводом, score matching, SDE и ODE.
\end{itemize}

\paragraph{Ограничения.}
\begin{itemize}
    \item Последовательная многошаговая генерация.
    \item Высокая стоимость многократного вызова денойзера.
    \item Чувствительность к расписанию шума и численному решателю.
    \item Компромисс между скоростью, разнообразием и качеством.
    \item В условных моделях чрезмерно сильное guidance может создавать
    артефакты.
    \item В латентных моделях качество ограничено автоэнкодером.
\end{itemize}

% ================================================================
\section{Связь трансформеров и диффузионных моделей}
% ================================================================

\subsection{Принципиальное различие}

Трансформер определяет способ вычисления представлений:
\[
    \text{входная последовательность}
    \longmapsto
    \text{выходная последовательность}.
\]
Диффузионная модель определяет вероятностный процесс генерации:
\[
    \text{шум}
    \longmapsto
    \text{данные через последовательное удаление шума}.
\]

Поэтому:
\begin{itemize}
    \item трансформер может быть недиффузионным;
    \item диффузионная модель может не содержать трансформера;
    \item трансформер может использоваться как денойзер в
    диффузионной модели;
    \item текстовый трансформер может кодировать условие для
    диффузионного генератора.
\end{itemize}

\subsection{Сравнительная таблица}

\begin{center}
\begin{tabular}{p{0.22\textwidth}p{0.35\textwidth}p{0.35\textwidth}}
\toprule
\textbf{Свойство} &
\textbf{Авторегрессионный трансформер} &
\textbf{Диффузионная модель}\\
\midrule
Базовый объект &
Условное распределение следующего токена &
Обратный переход от более шумного состояния к менее шумному\\
\addlinespace
Факторизация &
\[
    p(\bm{w})
    =
    \prod_t p(w_t\mid w_{<t})
\]
&
\[
    p_\theta(\bm{x}_{0:T})
    =
    p(\bm{x}_T)
    \prod_t p_\theta(\bm{x}_{t-1}\mid\bm{x}_t)
\]\\
\addlinespace
Обучение &
Cross-entropy по токенам &
MSE шума, score-функции, $\bm{x}_0$ или $\bm{v}$\\
\addlinespace
Порядок генерации &
По одному токену &
По одному временному шагу для всего объекта\\
\addlinespace
Параллельность одного шага &
Обрабатывается контекст, но новый токен один &
На шаге одновременно обновляются все координаты объекта\\
\addlinespace
Типичные данные &
Текст и дискретные последовательности &
Изображения, аудио, видео, непрерывные и структурированные данные\\
\addlinespace
Условие &
Префикс, специальные токены, cross-attention &
Конкатенация, cross-attention, classifier guidance, CFG\\
\bottomrule
\end{tabular}
\end{center}

\subsection{Общая схема text-to-image модели}

Типичная схема условной генерации изображения по тексту:
\begin{enumerate}
    \item текст разбивается на токены;
    \item трансформерный кодировщик строит матрицу текстовых
    представлений;
    \item из стандартного нормального распределения генерируется
    начальный латентный шум;
    \item U-Net или Diffusion Transformer на каждом временном шаге
    использует:
    \begin{itemize}
        \item текущий зашумлённый латентный код;
        \item эмбеддинг времени;
        \item текстовые признаки через cross-attention;
    \end{itemize}
    \item после удаления шума латентный код декодируется в изображение.
\end{enumerate}

\begin{figure}[h]
\centering
\begin{tikzpicture}[
    box/.style={
        draw,
        rounded corners,
        minimum width=3.5cm,
        minimum height=0.85cm,
        align=center,
        fill=blue!6
    },
    arrow/.style={-{Latex[length=2.5mm]},thick},
    node distance=0.75cm and 1.2cm
]
\node[box] (text) {Текст};
\node[box,right=of text] (tokens) {Токены};
\node[box,right=of tokens] (encoder) {Трансформерный\\кодировщик};

\node[box,below=1.5cm of tokens] (noise) {Латентный шум $\bm{z}_T$};
\node[box,right=of noise] (denoiser) {U-Net или DiT\\с cross-attention};
\node[box,right=of denoiser] (latent) {Чистый латент\\$\bm{z}_0$};
\node[box,right=of latent] (decoder) {Декодер};
\node[box,right=of decoder] (image) {Изображение};

\draw[arrow] (text)--(tokens);
\draw[arrow] (tokens)--(encoder);
\draw[arrow] (noise)--(denoiser);
\draw[arrow] (denoiser)--(latent);
\draw[arrow] (latent)--(decoder);
\draw[arrow] (decoder)--(image);
\draw[arrow] (encoder.south) |- (denoiser.north);
\end{tikzpicture}
\caption{Совместное использование трансформера и латентной
диффузионной модели.}
\end{figure}

% ================================================================
\section{Краткая формульная памятка}
% ================================================================

\subsection*{Трансформер}

\[
    \bm{Q}=\bm{X}_Q\bm{W}^Q,\qquad
    \bm{K}=\bm{X}_{KV}\bm{W}^K,\qquad
    \bm{V}=\bm{X}_{KV}\bm{W}^V.
\]

\[
    \operatorname{Attention}(\bm{Q},\bm{K},\bm{V})
    =
    \softmax\left(
        \frac{\bm{Q}\bm{K}^{\top}}{\sqrt{d_k}}+\bm{M}
    \right)\bm{V}.
\]

\[
    \operatorname{MHA}
    =
    \operatorname{Concat}
    (\operatorname{head}_1,\ldots,\operatorname{head}_h)\bm{W}^O.
\]

\[
    \operatorname{FFN}(\bm{x})
    =
    \phi(\bm{x}\bm{W}_1+\bm{b}_1)\bm{W}_2+\bm{b}_2.
\]

\[
    \bm{U}
    =
    \bm{X}
    +
    \operatorname{MHA}(\operatorname{LN}(\bm{X})),
\]
\[
    \bm{Y}
    =
    \bm{U}
    +
    \operatorname{FFN}(\operatorname{LN}(\bm{U})).
\]

\[
    p_\theta(w_{1:n})
    =
    \prod_{t=1}^{n}p_\theta(w_t\mid w_{<t}).
\]

\[
    \mathcal{L}_{\text{AR}}
    =
    -\sum_t\log p_\theta(w_t\mid w_{<t}).
\]

\subsection*{DDPM}

\[
    \alpha_t=1-\beta_t,
    \qquad
    \overline{\alpha}_t=\prod_{s=1}^{t}\alpha_s.
\]

\[
    q(\bm{x}_t\mid\bm{x}_{t-1})
    =
    \N(
        \sqrt{\alpha_t}\bm{x}_{t-1},
        \beta_t\bm{I}
    ).
\]

\[
    \bm{x}_t
    =
    \sqrt{\overline{\alpha}_t}\bm{x}_0
    +
    \sqrt{1-\overline{\alpha}_t}\bm{\eps}.
\]

\[
    \widetilde{\beta}_t
    =
    \frac{1-\overline{\alpha}_{t-1}}
         {1-\overline{\alpha}_t}\beta_t.
\]

\[
    \widetilde{\bm{\mu}}_t
    =
    \frac{
        \sqrt{\overline{\alpha}_{t-1}}\beta_t
    }{
        1-\overline{\alpha}_t
    }\bm{x}_0
    +
    \frac{
        \sqrt{\alpha_t}
        (1-\overline{\alpha}_{t-1})
    }{
        1-\overline{\alpha}_t
    }\bm{x}_t.
\]

\[
    \widehat{\bm{x}}_0
    =
    \frac{
        \bm{x}_t
        -
        \sqrt{1-\overline{\alpha}_t}
        \bm{\eps}_\theta(\bm{x}_t,t)
    }{
        \sqrt{\overline{\alpha}_t}
    }.
\]

\[
    \bm{\mu}_\theta
    =
    \frac{1}{\sqrt{\alpha_t}}
    \left(
        \bm{x}_t
        -
        \frac{\beta_t}
             {\sqrt{1-\overline{\alpha}_t}}
        \bm{\eps}_\theta(\bm{x}_t,t)
    \right).
\]

\[
    \mathcal{L}_{\text{simple}}
    =
    \E
    \left[
        \|\bm{\eps}
        -
        \bm{\eps}_\theta(\bm{x}_t,t)\|^2
    \right].
\]

\[
    \bm{s}_\theta(\bm{x}_t,t)
    =
    -
    \frac{
        \bm{\eps}_\theta(\bm{x}_t,t)
    }{
        \sqrt{1-\overline{\alpha}_t}
    }.
\]

\[
    \bm{\eps}_{\text{CFG}}
    =
    \bm{\eps}_{\text{uncond}}
    +
    w
    (\bm{\eps}_{\text{cond}}
    -
    \bm{\eps}_{\text{uncond}}).
\]

\subsection*{Непрерывная диффузия}

\[
    \dd\bm{X}_t
    =
    \bm{f}(\bm{X}_t,t)\dd t
    +
    g(t)\dd\bm{W}_t.
\]

Обратное SDE:
\[
    \dd\bm{X}_t
    =
    \left[
        \bm{f}(\bm{X}_t,t)
        -
        g(t)^2\nabla_{\bm{x}}\log p_t(\bm{X}_t)
    \right]\dd t
    +
    g(t)\dd\overline{\bm{W}}_t,
\]
где время идёт от $T$ к $0$.

Probability flow ODE:
\[
    \frac{\dd\bm{X}_t}{\dd t}
    =
    \bm{f}(\bm{X}_t,t)
    -
    \frac{1}{2}g(t)^2
    \nabla_{\bm{x}}\log p_t(\bm{X}_t).
\]

\end{document}
